\subsection{Dimensión relativa de un preesquema liso sobre otro}
\label{subsection:IV.17.10}

\begin{definition}[17.10.1]
\label{IV.17.10.1}
Sea $f:X\to Y$ un morfismo localmente de tipo finito. Se llama \emph{dimensión relativa de $f$ en el punto $x\in X$} (o \emph{dimensión relativa de $X$ sobre $Y$ en el punto $x$}) y se denota por $\dim_x f$ el entero no negativo $\dim_x(f^{-1}(f(x)))$.
\end{definition}

Decir que $f$ es \emph{cuasifinito} en el punto $x$ (Err$_{\mathrm{III}}$, 20) equivale, por tanto, a decir que $\dim_x f=0$. Hemos visto (13.1.3) que la función $x\mapsto\dim_x f$ es semicontinua superiormente. Obsérvese que, incluso cuando el morfismo $f$ posee la propiedad $(S_1)$ (es decir, por (6.8.1), es plano y sus fibras no tienen ningún ciclo primo asociado inmerso), la función $x\mapsto\dim_x f$ no es necesariamente continua, como muestra el ejemplo en que $Y=\Spec(k)$, donde $k$ es un cuerpo, y $X=\Spec(k[U,V,W]/\mathfrak p\mathfrak q)$, donde $\mathfrak p=(W)$ y $\mathfrak q=(U)+(V-W)$ son ideales primos de $k[U,V,W]$ (siendo, por tanto, $X$ la unión, en el espacio de 3 dimensiones, de un plano y una recta no paralela a dicho plano).

Decir que un morfismo localmente de presentación finita $f:X\to Y$ es étale en el punto $x\in X$ significa asimismo que $f$ es liso en el punto $x$ y que $\dim_x f=0$ (17.6.1).

\begin{proposition}[17.10.2]
\label{IV.17.10.2}
Sea $f:X\to Y$ un morfismo liso. Para todo $x\in X$, el $\sh{O}_X$-módulo localmente libre $\Omega_f^1$ (17.2.3) tiene en $x$ un rango igual a $\dim_x f$ (lo que implica que $x\mapsto\dim_x f$ es una función continua en $X$).
\end{proposition}

En efecto, si $y=f(x)$, $X_y=f^{-1}(y)$ es liso sobre $k(y)$, y si $f_y:X_y\to\Spec(k(y))$ es el morfismo estructural, se tiene $\Omega_{f_y}^1=\Omega_f^1\otimes_{\sh{O}_X}\sh{O}_{X_y}$ (16.4.5); podemos, por tanto, limitarnos al caso en que $Y=\Spec(k)$ es el espectro de un cuerpo; además, en virtud de (16.4.5) y (17.7.1), podemos reemplazar $k$ por una extensión algebraicamente cerrada, dicho de otro modo,

\oldpage[IV-4]{82}

suponer que $k$ es algebraicamente cerrado. Como el conjunto de puntos de $X$ racionales sobre $k$ es entonces denso en $X$ (10.4.8), podemos limitarnos al caso en que $x$ es racional sobre $k$. Ahora bien, (16.4.12) $(\Omega^1_{X/k})_x\otimes_{\sh{O}_{X,x}}k(x)$ es entonces $k$-isomorfo a $\mathfrak m_x/\mathfrak m_x^2$, y puesto que $\sh{O}_x$ es un anillo local regular (17.5.1), $\operatorname{rg}_{k(x)}(\mathfrak m_x/\mathfrak m_x^2)=\dim(\sh{O}_x)$ (0, 17.1.1); por tanto, $\operatorname{rg}_{\sh{O}_x}(\Omega^1_{X/k})_x=\dim(\sh{O}_x)$. Pero, puesto que $k(x)=k$, se tiene $\dim_x f=\dim(\sh{O}_x)$ (5.2.3). C.Q.D.

Más adelante demostraremos un recíproco de este resultado (17.15.5).

\begin{corollary}[17.10.3]
\label{IV.17.10.3}
Sean $f:X\to Y$, $g:Y\to Z$ dos morfismos lisos. Entonces, para todo $x\in X$, se tiene
\begin{equation}
\label{IV.17.10.3.1}
\tag{17.10.3.1}
\dim_x(g\circ f)=\dim_x f+\dim_{f(x)}g.
\end{equation}
\end{corollary}

En efecto, $g\circ f$ es liso (17.3.3), por lo que los tres $\sh{O}_X$-módulos $\Omega^1_{X/Y}$, $\Omega^1_{X/Z}$ y $f^*(\Omega^1_{Y/Z})$ son localmente libres (17.2.3 y $0_{\mathrm I}$, 5.4.5); además, el rango en $x$ de $f^*(\Omega^1_{Y/Z})$ es igual al rango en $y=f(x)$ de $\Omega^1_{Y/Z}$. La igualdad (17.10.3.1) es, por tanto, consecuencia de (17.10.2) y de la exactitud de la sucesión (17.2.3.1).

\begin{corollary}[17.10.4]
\label{IV.17.10.4}
Sean $f:X\to Y$ un morfismo liso y $X'$ un subpreesquema de $X$ tal que el morfismo compuesto $X'\xrightarrow{j}X\xrightarrow{f}Y$ (donde $j$ es la inyección canónica) sea liso. Entonces el haz conormal $\sh{N}_{X'/X}$ es un $\sh{O}_{X'}$-módulo localmente libre, y para todo $x\in X'$, se tiene
\begin{equation}
\label{IV.17.10.4.1}
\tag{17.10.4.1}
\dim_x f=\dim_x(f\circ j)+\operatorname{rg}_{\sh{O}_{X',x}}(\sh{N}_{X'/X})_x.
\end{equation}
\end{corollary}

En efecto, $\Omega^1_{X/Y}\otimes_{\sh{O}_X}\sh{O}_{X'}$ y $\Omega^1_{X'/Y}$ son entonces localmente libres y la sucesión exacta (17.2.5.1) se escinde en una vecindad adecuada de cada punto de $X'$, por lo que $\sh{N}_{X'/X}$ es localmente libre (Bourbaki, \emph{Alg. comm.}, cap. II, § 5, n\textsuperscript{o} 2, teorema 1), y la relación (17.10.4.1) resulta inmediatamente de la exactitud de la sucesión (17.2.5.1).
